\section{Discussion}
\label{sec:discussion}

\subsection{Interpretation of Findings}

\para{Semantic matching is necessary but not sufficient.}
The 40\% top-1 accuracy of cosine similarity confirms that semantic content is a genuine
driver of tool selection, but the substantial gap to the LLM's actual accuracy (and the
gap between selected vs.\ top-similar tool cosim) shows that the model's internal semantic
matching is more discriminating than simple embedding distance.  The LLM likely incorporates
joint reasoning over query intent and tool capability that embedding similarity
cannot capture.  This is consistent with the BiasBusters finding \citep{blankenstein2025biasbust}
that semantic alignment is the ``strongest predictor'' while not being deterministic.

\para{The category vs.\ position dissociation reveals a processing hierarchy.}
The most striking mechanistic finding is the clean separation between category encoding
and position selection in \gpttwo.  Category is fully encoded from layer~1---suggesting
it is a largely lexical property, encoded in the token embedding or the first
self-attention layer, and does not require deep processing.  This aligns with the
``E-step hypothesis'' in the MI literature \citep{sharkey2025open}: early layers handle
lexical and syntactic features, while later layers handle compositional and relational
computations.  Within-category tool identification, by contrast, requires comparing the
query against each tool description---a relational computation that reaches maximum
discriminability only at layer~9.  The subsequent drop at layer~12 may reflect the final
layer specializing for next-token prediction rather than maintaining a discriminative
comparison representation.

\para{Tool name as a secondary selection signal.}
The generic-description experiment (no change to selection) and the translation
adversarial case (no change despite description swap) both reveal that tool \emph{name}
carries its own strong prior.  ``Google Translate'' is a distinctive, highly familiar
name; its prior overwhelms the adversarial description.  In contrast, ``WolframAlpha'' and
``OpenWeatherMap'' are less unambiguously associated with their categories in natural
language, so description content dominates for those tools.  This suggests a two-path
selection mechanism: a fast name-based prior and a slower description-based semantic
matching process, similar to dual-process models in cognitive science.

\para{Last-token as selection buffer.}
The high and consistent last-token probe accuracy (73--81\% across all layers) versus the
gradual rise in mean-pooling accuracy suggests that the last-token representation
functions as an active selection buffer.  In autoregressive generation, the last token
must carry the information needed to generate the tool name; the residual stream at that
position aggregates relevant context from the very first layer.  This is mechanistically
consistent with the ``copy mechanism'' in IOI circuits \citep{merullo2024circuit}, where
the final token representation aggregates identity information about the target entity.

\subsection{Limitations}

\para{Small sample size.}
The main behavioral analysis uses 45 LLM selections (15 per run $\times$ 3 runs), which
limits statistical power.  Per-category analyses have at most $n = 9$ examples;
the chi-square test for positional bias has 40 observations.  Larger datasets ($n \geq
500$) would enable Bonferroni-corrected per-category comparisons and tighter confidence
intervals.

\para{GPT-2-small is not an instruction-tuned tool-calling model.}
Our mechanistic findings come from \gpttwo, a base language model without tool-use
training.  The representations we probe reflect general language model processing, not
specialized tool-selection circuits.  Instruction-tuned models fine-tuned on function
calling (e.g., LLaMA-3.1-8B-Instruct, GPT-4.1-mini) may develop qualitatively different
representations or dedicate specific heads to tool routing.  The \citet{tigges2024llm}
finding that circuit structures are consistent across model scale provides \emph{some}
justification for the \gpttwo results, but direct probing of larger instruction-tuned
models is essential future work.

\para{Probing detects correlation, not causality.}
Linear probes measure the decodability of information from representations, but high
probe accuracy does not prove that the probed information is causally used for the
behavior \citep{sharkey2025open}.  The perturbation experiments provide complementary
causal evidence at the behavioral level, but activation patching on specific attention
heads---the method used in circuit-level ACDC analysis \citep{conmy2023acdc}---would be
needed to establish causal responsibility at the circuit level.

\para{No circuit-level analysis.}
This work identifies \emph{where} information is encoded but not \emph{which} circuits
compute it.  ACDC circuit discovery \citep{conmy2023acdc} was not applied; applying it
to \gpttwo tool-selection prompts would require GPU resources and would identify the
specific attention heads performing cross-tool comparison in layers 7--9.

\para{Fixed tool descriptions.}
Our curated descriptions are drawn from real API documentation but are fixed in length
and style.  Real-world tool registries vary enormously in description quality.  Systems
with short, vague, or redundant descriptions would likely show lower semantic accuracy
and higher position bias than reported here.

\subsection{Broader Implications}

These findings have direct engineering implications.  First, within-category tool
selection is the hardest case (37.5\% semantic accuracy)---yet it is the most common
scenario in multi-tool agents that register many tools with overlapping functions.
Practitioners should prioritize making within-category descriptions maximally distinctive.
Second, tool names carry independent priors: agents deploying tools with highly
recognizable names (e.g., Google-branded APIs) may find that description content is
partially overridden.  Third, the layer~9 peak in position probing suggests that
interpretability interventions---such as activation steering---applied at that layer
could potentially redirect tool selection without modifying the model's weights, an
exciting direction for robust tool-routing systems.
